\documentclass[11pt]{article}
\usepackage[margin=1in]{geometry}
\usepackage{amsmath,amssymb,amsthm}
\usepackage[hidelinks]{hyperref}
\newtheorem{theorem}{Theorem}
\title{Nonconvexity of non-completable probability tables\\
\large A counterexample to MathDB 400968}
\author{Principia Math in Progress}
\date{12 September 2026}
\begin{document}
\maketitle
\section*{Task 1. Resolve the original convexity question}
The problem in Kummer's Oberwolfach report \cite[p.~1271]{report}
asks whether the nonnegative partial tensors which have no completion
as the joint distribution of independent random variables always form
a convex set. The stated setting has a full-dimensional completion
image and $|E|=\sum_i(d_i-1)$ observed coordinates.

\begin{theorem}
The answer is negative, already for a $3\times3$ table with
\[
 E=\{(1,1),(1,2),(2,1),(3,3)\}.
\]
This pattern observes every row and every column, and its completion
image has nonempty interior in $\mathbb R^4$.
\end{theorem}
\begin{proof}
Use the coordinate order $(11,12,21,33)$, and define
\[
 P=\left(\frac1{32},\frac18,\frac18,\frac1{16}\right),
 \qquad
 Q=\left(\frac{15}{32},\frac18,\frac18,\frac1{16}\right).
\]
All entries are positive, and their respective observed sums are
$11/32$ and $25/32$, both less than one.

We first obtain a necessary condition for completion. Write a putative
completed table as $T_{ij}=a_i b_j$, where $a,b$ are nonnegative
probability vectors. If its observed coordinates are $(x,y,z,w)$
with $x>0$, rank one forces $T_{22}=yz/x$. Consequently
\[
 S:=x+y+z+\frac{yz}{x}=(a_1+a_2)(b_1+b_2),
 \qquad w=a_3b_3.
\]
Cauchy--Schwarz, applied to the two unit vectors
$(\sqrt{1-a_3},\sqrt{a_3})$ and
$(\sqrt{1-b_3},\sqrt{b_3})$, gives
\[
 \sqrt S+\sqrt w\leq1.
\]
In particular, when $w=1/16$, every completion must satisfy
$S\leq9/16$. For our two points, however,
\[
 S(P)=\frac{25}{32}>\frac9{16},
 \qquad
 S(Q)=\frac{361}{480}>\frac9{16}.
\]
Thus both $P$ and $Q$ are non-completable.

Their midpoint is completable:
\[
 \frac{P+Q}{2}
 =\left(\frac14,\frac18,\frac18,\frac1{16}\right),
 \qquad
 T=\frac1{16}
 \begin{pmatrix}
 4&2&2\\2&1&1\\2&1&1
 \end{pmatrix}
 =
 \begin{pmatrix}1/2\\1/4\\1/4\end{pmatrix}
 \begin{pmatrix}1/2&1/4&1/4\end{pmatrix}.
\]
The entries of $T$ are positive and sum to one, and its entries in
$E$ are exactly the midpoint. This proves nonconvexity.

It remains to check the full-dimensional hypothesis.
Parameterize the two probability vectors by
$a=(a_1,a_2,1-a_1-a_2)$ and
$b=(b_1,b_2,1-b_1-b_2)$. The observed-coordinate map is
\[
 F(a_1,a_2,b_1,b_2)=
 \bigl(a_1b_1,a_1b_2,a_2b_1,
       (1-a_1-a_2)(1-b_1-b_2)\bigr).
\]
Its Jacobian determinant is
\[
 \det DF=a_1b_1(a_1+a_2-b_1-b_2).
\]
At the interior parameter point
$(a_1,a_2,b_1,b_2)=(1/2,1/4,1/3,1/3)$ it equals $1/72$.
The inverse function theorem therefore shows that the image contains
an open subset of $\mathbb R^4$. Finally,
$|E|=4=(3-1)+(3-1)$, as required.
\end{proof}

The necessary completion inequality used above agrees with the
general matrix-completion framework of Kubjas and Rosen \cite{KR}.
The counterexample is verified directly here, without requiring a
sufficiency criterion for arbitrary partial tables.

\section*{Task 2. Classify the convex matrix patterns}
Characterize the observed bipartite forests $E$ with
$|E|=d_1+d_2-2$ for which the set of non-completable positive partial
tables is convex. Include a criterion stated in terms of the connected
components of the observation graph.

\section*{Task 3. Determine the smallest counterexample patterns}
Determine the smallest table formats and observation counts admitting
nonconvexity when every row and column is observed and the completion
image has nonempty interior.

\section*{Task 4. Describe the convex hull}
For the displayed $3\times3$ pattern, determine the convex hull of the
non-completable region inside
$\{(x,y,z,w)\geq0:x+y+z+w\leq1\}$, including its boundary faces.

\section*{Task 5. Extend the analysis to higher-order tensors}
Find structural criteria for convexity or nonconvexity for tensor
orders at least three, requiring every coordinate value to occur in
the observation pattern and retaining the full-dimensional hypothesis.

\begin{thebibliography}{9}
\bibitem{report}
M.~Kummer, Rank One Tensor Completion, in \emph{Algebraic Statistics},
Oberwolfach Reports 14 (2017), 1270--1272, published 2018.
\href{https://doi.org/10.4171/OWR/2017/20}{doi:10.4171/OWR/2017/20}.
\bibitem{KR}
K.~Kubjas and Z.~Rosen, Matrix Completion for the Independence Model,
Journal of Algebraic Statistics 8 (2017), 1--21.
\href{https://arxiv.org/abs/1407.3254}{arXiv:1407.3254}.
\end{thebibliography}
\end{document}
